\clearpage
\VerdierPhysicalPageAnchor{0192}%
\begingroup\small\setlength{\parskip}{0pt}\setlength{\emergencystretch}{1.5em}
\setlength{\abovedisplayskip}{4pt plus 2pt minus 2pt}%
\setlength{\belowdisplayskip}{4pt plus 2pt minus 2pt}%

\VerdierSemanticAnchor{ch3d-2-2-5-cont}
整数 $n$ で $M\text{-}\dim(X)\leq n$ となるものが存在するとき、対象 $X$ の $M$-次元
（$M\text{-}\dim(X)$ と記す）とは、最小の整数 $n$ であって
$M\text{-}\dim(X)\leq n$ を満たすものをいう。そのような整数が存在しないときは
$M\text{-}\dim(X)=+\infty$ と定める。

\VerdierSemanticAnchor{ch3d-2-2-6}
\medskip\noindent\textbf{補題 2.2.6.}\quad
\textit{$0\longrightarrow X\longrightarrow L\xrightarrow{u}Y\longrightarrow0$
を $\mathcal A$ における完全列とし、$L\in M$ かつ $X,Y\in M'$ とする。このとき}
\[
M\text{-}\dim(X)=\sup\{M\text{-}\dim(Y)-1,0\}.
\]

\VerdierSemanticAnchor{ch3d-2-2-7}
\medskip\noindent\textbf{命題 2.2.7.}\quad
\textit{$0\longrightarrow X\longrightarrow Y\longrightarrow Z\longrightarrow0$
を完全列とし、$X,Y,Z\in M'$ とする。}

\noindent\textit{a) 
$M\text{-}\dim(Z)\leq
\inf\{M\text{-}\dim(X),M\text{-}\dim(Y)\}$ ならば、
$M\text{-}\dim(X)=M\text{-}\dim(Y)$ である。}

\noindent\textit{b) 
$M\text{-}\dim(Y)\leq
\inf\{M\text{-}\dim(X),M\text{-}\dim(Z)\}$ ならば、}
\[
M\text{-}\dim(X)=
\sup\{M\text{-}\dim(Z)-1,M\text{-}\dim(Y)\}.
\]

\noindent\textit{c) 
$M\text{-}\dim(X)\leq
\inf\{M\text{-}\dim(Y),M\text{-}\dim(Z)\}$ ならば、
$M\text{-}\dim(Z)=M\text{-}\dim(Y)$ または
$M\text{-}\dim(Z)=M\text{-}\dim(Y)+1$ である。}

\VerdierSemanticAnchor{ch3d-2-2-7-proof}
完全列の間のエピ射を構成し、帰納法を用いて命題を証明する：

\VerdierSemanticAnchor{ch3d-2-2-7-epimorphism-diagram}
\begin{center}
\begin{tikzcd}[column sep=2.8em,row sep=2.4em]
0 \arrow[r] & L_2 \arrow[r] \arrow[d] & L_1 \arrow[r] \arrow[d]
  & L_0 \arrow[r] \arrow[d] & 0 \\
0 \arrow[r] & X \arrow[r] & Y \arrow[r] & Z \arrow[r] & 0
\end{tikzcd}
\end{center}
\[
L_i\in M,\qquad 0\leq i\leq2,
\]
そして、補題 (2.2.6) によって制御される核の完全列に帰納法の仮定を適用する。
したがって、補題を証明すればよい。

\VerdierSemanticAnchor{ch3d-2-2-6-proof}
まず、$M\text{-}\dim(Y)=+\infty$（それぞれ $0$）ならば、
$M\text{-}\dim(X)=+\infty$（それぞれ $0$）である。したがって、
$M\text{-}\dim(Y)$ は有限かつ零でないと仮定してよく、証明すべきことは
$M\text{-}\dim(X)=M\text{-}\dim(Y)-1$ である。次の

\VerdierSemanticAnchor{ch3d-2-2-6-resolution}
\[
0\longrightarrow L_n\longrightarrow L_{n-1}\longrightarrow\cdots
\longrightarrow L_0\xrightarrow{u}Y\longrightarrow0,
\qquad n\geq1,
\]
を $Y$ の $M$-分解とする。次をおく：

\VerdierSemanticAnchor{ch3d-2-2-6-v-complex}
\[
\begin{aligned}
V^\bullet={}&0\longrightarrow0\longrightarrow\cdots\longrightarrow0
\longrightarrow L_n\longrightarrow L_{n-1}\longrightarrow\cdots\\
&\longrightarrow L_0\xrightarrow{u}Y\longrightarrow0
\longrightarrow0\longrightarrow\cdots
\end{aligned}
\]
（$L_i$ は次数 $-i$ の成分である）。さらに、

\VerdierSemanticAnchor{ch3d-2-2-6-u0}
\[
U_0^\bullet=V^\bullet(u,1),
\qquad
p_0=p(u,1)\colon U_0^\bullet\longrightarrow V^\bullet.
\]
\endgroup
